%%%%%%%%%%%%%%%%%%%%%%%%%%%%%%%%%%%%%%%%%
% Academic Paper Response Letter. 
% LaTeX Template
% Version 1.0 (20/07/2023)
% 
% This template refers to the toolbox tutorial (in Chinese):
% https://liam.page/2016/07/22/using-the-tcolorbox-package-to-create-a-new-theorem-environment/
% Please refer to the tutorial for your customization. 
% 
% Template license: Creative Commons CC BY 4.0
%
% Author: Yuhan Xie
% Contact email: yhxie0107@gmail.com
% Any revision advice is welcome!
% 
%%%%%%%%%%%%%%%%%%%%%%%%%%%%%%%%%%%%%%%%%


\documentclass{article}
%% Language and font encodings
\usepackage[english]{babel}
\usepackage{times} % make the section title in Times New Roman

\usepackage{booktabs}
\usepackage{tabu}
\usepackage[T1]{fontenc}

%% Sets page size and margins
\usepackage[a4paper,top=2.5cm,bottom=2cm,left=1.7cm,right=1.7cm,marginparwidth=1.75cm]{geometry}

%% Useful packages
\usepackage{amsmath}
\usepackage{graphicx}
\usepackage[colorinlistoftodos]{todonotes}
\usepackage[colorlinks=true, allcolors=blue]{hyperref}
\usepackage{indentfirst}
\usepackage{graphicx}
\usepackage{subfigure}
\usepackage{float}
\usepackage{threeparttable} 
\usepackage{multirow}
\usepackage{url}


\usepackage[most]{tcolorbox}


%% Comment box environment used throughout the letter
\newtcbtheorem[auto counter]{cmt}{Comment}{
	colbacktitle = black!60!white, colframe = black!60!white,
	colback = black!5!white,
	fonttitle=\bfseries,%fontupper=\itshape,
}{t}


\begin{document}

%%%%%%%%%%%%%%%%%%%%%%%%%%%%%%%%%%%%%%%%%%%%%%%%%%%%%%%%%%%%%%%%%%%%%%%%%%%%%%%
% Response to Reviewers -- PONE-D-26-28965
%
% Structure follows POINT_BREAKDOWN.md (preliminary-structure version):
% four sections, no separate "Response to the Editor" section --- the
% journal-side items are grouped under "Journal Requirements".
%   1. Journal Requirements ........  7 items
%   2. Reviewer #1 ................. 11 items
%   3. Reviewer #2 ................. 14 items
%   4. Reviewer #3 ................. 12 items
%   Total ......................... 44 response items
%
% Comment boxes quote the review verbatim. The state of each response is marked
% visibly in the body of the letter (see the banner on the first page), never
% here: %-comments do not reach the PDF that the authors actually read.
%%%%%%%%%%%%%%%%%%%%%%%%%%%%%%%%%%%%%%%%%%%%%%%%%%%%%%%%%%%%%%%%%%%%%%%%%%%%%%%


\begin{center}
{\Large\bfseries Response to Reviewers}

\vspace{0.2cm}
Manuscript PONE-D-26-28965

\vspace{0.1cm}
\textit{Evaluation of validation strategies and transferability of EEG attention
decoding: a comparative analysis on open datasets}
\end{center}

\vspace{0.3cm}

\noindent
\textbf{[TODO --- working state of this letter; this banner is not part of the
submission.]} Every response below is in one of three states. A response with no
marker is backed both by the code in \texttt{python\_project/} and by text that is
actually present in \texttt{manuscript\_clean/manuscript.tex}. A response
followed by a \textbf{[TODO]} paragraph is only partly backed: the
implementation exists, but the manuscript does not yet say so, or the numbers it
refers to have not been computed. A response consisting only of \textbf{[TODO]}
is not answered at all yet. The manuscript is mid-rewrite --- its Abstract,
Background, Models and Results sections are cleared placeholders --- the full
sweep in \texttt{python\_project/results/} is still running, and no analysis has
been produced from it. Until that changes, no response here may cite a number, a
table, a figure or a manuscript passage. All markers must be gone before the
letter is uploaded.

\vspace{0.3cm}
\noindent
Dear Editor and Reviewers,

\vspace{0.2cm}
\noindent
We thank the Editor and the Reviewers for their careful evaluation of our
manuscript and for their constructive comments. We have carefully considered all
comments and are preparing a revised version of the manuscript that addresses
the issues raised during peer review.

\vspace{0.2cm}
\noindent
The major revisions will include:

\begin{itemize}
	\item \textbf{[TODO --- major revision 1: redefinition of the target
	      construct.]} The direction is settled --- the study claims
	      discrimination between operationally defined experimental conditions
	      rather than decoding of attention as a psychological construct --- but
	      the manuscript does not yet carry it: Table~1 still names its classes
	      ``High Attention'' and ``Low Attention'', and no label-mapping
	      sensitivity analysis exists in the code. Only the exclusion of drowsy
	      epochs from the primary analysis is in place. Rewrite this bullet
	      together with the responses to MR1 and Comments~3.1, 3.2 and 4.5.
	\item \emph{[Major revision 2 --- to be completed after the corresponding manuscript changes are finalized.]}
	\item \emph{[Major revision 3 --- to be completed after the corresponding manuscript changes are finalized.]}
	\item \emph{[Additional major revision, if needed.]}
\end{itemize}

\noindent
For clarity, our response is organized into four sections:

\begin{enumerate}
	\item \textbf{Journal Requirements}
	\item \textbf{Reviewer \#1}
	\item \textbf{Reviewer \#2}
	\item \textbf{Reviewer \#3}
\end{enumerate}

\noindent
Reviewer and journal comments are reproduced below and followed by our
point-by-point responses. Where a single original paragraph contains several
independent criticisms or requests, it is divided into separate response points
so that each issue can be addressed explicitly. Changes to the manuscript are
identified with the corresponding page and line numbers wherever applicable.

\vspace{0.4cm}
\noindent
Yours sincerely,

\vspace{0.2cm}
\noindent
\emph{[TODO: author list]}


%%%%%%%%%%%%%%%%%%%%%%%%%%%%%%%%%%%%%%%%%%%%%%%%%%%%%%%%%% Journal Requirements
\newpage
\setcounter{tcb@cnt@cmt}{0}
\section*{Journal Requirements}

\noindent
The seven requirements listed in the decision letter are addressed below.

\vspace{0.3cm}

%%%%%%%%%%%%% Journal Requirement 1
\begin{cmt}{}{}

Please ensure that your manuscript meets PLOS ONE's style requirements,
including those for file naming. The PLOS ONE style templates can be found at
\url{https://journals.plos.org/plosone/s/file?id=wjVg/PLOSOne_formatting_sample_main_body.pdf}
and
\url{https://journals.plos.org/plosone/s/file?id=ba62/PLOSOne_formatting_sample_title_authors_affiliations.pdf}

\end{cmt}
\vspace{0.1cm}
\noindent
\underline{\textbf{Response:}}
\vspace{0.2cm}

\emph{[TODO]}

\vspace{0.3cm}


%%%%%%%%%%%%% Journal Requirement 2
\begin{cmt}{}{}

Please note that PLOS One has specific guidelines on code sharing for
submissions in which author-generated code underpins the findings in the
manuscript. In these cases, we expect all author-generated code to be made
available without restrictions upon publication of the work. Please review our
guidelines at
\url{https://journals.plos.org/plosone/s/materials-and-software-sharing#loc-sharing-code}
and ensure that your code is shared in a way that follows best practice and
facilitates reproducibility and reuse.

\end{cmt}
\vspace{0.1cm}
\noindent
\underline{\textbf{Response:}}
\vspace{0.2cm}

\emph{[TODO]}

\vspace{0.3cm}

%%%%%%%%%%%%% Journal Requirement 3
\begin{cmt}{}{}

We note that the grant information you provided in the `Funding Information'
and `Financial Disclosure' sections do not match.
When you resubmit, please ensure that you provide the correct grant numbers for
the awards you received for your study in the `Funding Information' section.

\end{cmt}
\vspace{0.1cm}
\noindent
\underline{\textbf{Response:}}
\vspace{0.2cm}

\emph{[TODO]}

\vspace{0.3cm}


%%%%%%%%%%%%% Journal Requirement 4
\begin{cmt}{}{}

Thank you for stating the following financial disclosure:

``The author(s) declared that financial support was received for this work
and/or its publication. This research was funded by the Committee of Science of
the Ministry of Science and Higher Education of the Republic of Kazakhstan
(Grant no. BR27198099).''

Please state what role the funders took in the study. If the funders had no
role, please state: ``The funders had no role in study design, data collection
and analysis, decision to publish, or preparation of the manuscript.''

If this statement is not correct you must amend it as needed.

Please include this amended Role of Funder statement in your cover letter; we
will change the online submission form on your behalf.

\end{cmt}
\vspace{0.1cm}
\noindent
\underline{\textbf{Response:}}
\vspace{0.2cm}

\emph{[TODO]}

\vspace{0.3cm}


%%%%%%%%%%%%% Journal Requirement 5
\begin{cmt}{}{}

Please note that your Data Availability Statement is currently missing the
repository name and/or the DOI/accession number of each dataset OR a direct link
to access each database. If your manuscript is accepted for publication, you
will be asked to provide these details on a very short timeline. We therefore
suggest that you provide this information now, though we will not hold up the
peer review process if you are unable.

\end{cmt}
\vspace{0.1cm}
\noindent
\underline{\textbf{Response:}}
\vspace{0.2cm}

\emph{[TODO]}

\vspace{0.3cm}


%%%%%%%%%%%%% Journal Requirement 6
\begin{cmt}{}{}

We notice that your supplementary figures are uploaded with the file type
`Figure'. Please amend the file type to `Supporting Information'. Please ensure
that each Supporting Information file has a legend listed in the manuscript
after the references list.

\end{cmt}
\vspace{0.1cm}
\noindent
\underline{\textbf{Response:}}
\vspace{0.2cm}

\emph{[TODO]}

\vspace{0.3cm}


%%%%%%%%%%%%% Journal Requirement 7
\begin{cmt}{}{}

If the reviewer comments include a recommendation to cite specific previously
published works, please review and evaluate these publications to determine
whether they are relevant and should be cited. There is no requirement to cite
these works unless the editor has indicated otherwise.

\end{cmt}
\vspace{0.1cm}
\noindent
\underline{\textbf{Response:}}
\vspace{0.2cm}

\emph{[TODO]}

\vspace{0.3cm}


%%%%%%%%%%%%%%%%%%%%%%%%%%%%%%%%%%%%%%%%%%%%%%%%%%%%%%%%%%%%%%% Reviewer 1
\newpage
\setcounter{tcb@cnt@cmt}{0}
\section*{Reviewer \#1}

\vspace{0.4cm}

%%%%%%%%%%%% Comment 2.1 -- MR1
\begin{cmt}{}{}

MR1. The binary attention labels have uncertain construct validity. Dataset B
was collected for stress analysis rather than attention measurement, while
Dataset A uses instructed states presented in a fixed temporal sequence.
Drowsiness also involves eye closure. Consequently, the classes confound
attention with stress, arousal, visual input, ocular activity and time-on-task.

\end{cmt}
\vspace{0.1cm}
\noindent
\underline{\textbf{Response:}}
\vspace{0.2cm}

We thank the Reviewer for identifying this fundamental issue. We agree that the
original framing overstated what the available labels can support: they are not
independent measurements of an individual's attention, and differences in
stress, arousal, sensory input, ocular activity, time-on-task, and other
protocol features may contribute to discriminability.

We therefore revised the research question, motivation, and knowledge gap
throughout the manuscript. The study no longer treats its results as evidence
for or against invariant neural markers of attention. Instead, it evaluates
zero-shot transfer of a \emph{protocol-defined} contrast between higher-demand
conditions and rest or task disengagement, and reports pre-adaptation benchmarks
for source--target suitability and negative transfer. The revised Introduction
states this boundary of inference explicitly. We also excluded the
\textit{Drowsy} phase because eye closure and drowsiness would add pronounced
differences in arousal, visual input, and ocular activity. The remaining
confounds are acknowledged as limitations rather than interpreted as attention
effects.

\vspace{0.3cm}


%%%%%%%%%%%% Comment 2.2 -- MR2
\begin{cmt}{}{}

MR2. The principal within-dataset baseline is affected by substantial
information leakage. One-second windows separated by only 0.125 seconds share
87.5\% of their samples, yet temporally adjacent windows are distributed across
stratified folds. Acknowledging this issue as a limitation does not make the
reported baseline scientifically valid.

\end{cmt}
\vspace{0.1cm}
\noindent
\underline{\textbf{Response:}}
\vspace{0.2cm}

We thank the Reviewer for drawing our attention to this issue. The original
random split did not adequately account for the dependence between densely
overlapping windows.

We reconstructed the baseline using group-aware splitting at the recording-unit
level: all windows from the same recording unit are assigned wholly to either
the training or the test partition. Thus, no temporally adjacent overlapping
windows are split between partitions. We also replaced the previous 1-s windows
with 5-s windows and 0.5-s overlap, and recalculated all reported results under
this leakage-resistant design.

\vspace{0.3cm}


%%%%%%%%%%%% Comment 2.3 -- MR3
\begin{cmt}{}{}

MR3. Several validation procedures are insufficiently defined. Dataset B
trials appear to be treated as sessions without justification; the
composition of the ``task-only'' datasets is unclear; class counts are
absent; and cross-task testing does not separate task transfer from
participant identity.

\end{cmt}
\vspace{0.1cm}
\noindent
\underline{\textbf{Response:}}
\vspace{0.2cm}

We thank the Reviewer for this clarification. We do not treat the separately
provided SAM-40 task recordings as session-like: the source publication
describes the three repetitions as trials and does not establish separate
physical sessions. Cross-session validation therefore applies only to
Dataset~A.

\emph{Dataset compositions and class counts.} We added
Table~\ref{tab:dataset_compositions}, which defines the full, task-only, and
two-task Dataset~B compositions by their retained conditions. It also reports
the number of participants, recording units, and prepared pre-oversampling
windows in the low- and high-attention classes for each composition.

\emph{Subject-identity confounding in cross-task testing.} We rebuilt
cross-task validation as leave-one-subject-out transfer: for each fold, the
source-task training data exclude subject $S$, and the target-task test data
contain only subject $S$. Source and target partitions are therefore
subject-disjoint by construction.

\vspace{0.3cm}


%%%%%%%%%%%% Comment 2.4 -- MR4
\begin{cmt}{}{}

MR4. Reproducibility is incomplete within the manuscript. Filter type and order,
ICA implementation and rejection criteria, STFT window and FFT parameters,
feature dimensionality, scaler-fitting scope, XGBoost hyperparameters, tuning
procedure, random seeds and fold construction are not adequately reported.

\end{cmt}
\vspace{0.1cm}
\noindent
\underline{\textbf{Response:}}
\vspace{0.2cm}

We thank the Reviewer for identifying the insufficient methodological detail in
the original submission. We revised the Methodology throughout, specifically in
the \emph{Data preparation}, \emph{Models}, and \emph{Validation strategies}
sections.

\emph{Data preparation.} We now specify the filtering sequence and parameters,
the extended-Infomax ICA procedure, its ICLabel-based artifact-rejection
criterion and random state, the common channel montage, and the 5-s windows
with 0.5-s overlap. The revised description specifies 50-Hz notch filtering,
1--45-Hz band-pass filtering, average rereferencing, and removal of
eye-blink or muscle components classified with probability at least 0.80.

\emph{Models.} We now report the fixed 108-dimensional classical feature
representation (12 channels $\times$ 9 features per channel), the
training-fold-only fitting of data-dependent transformations, the fixed random
seed (42), and the absence of data-driven hyperparameter tuning.

\emph{Validation strategies.} We now provide a fixed description of fold
construction for every validation protocol, including the recording-unit
grouping that keeps all overlapping windows on one side of a training--test
split.

\vspace{0.3cm}


%%%%%%%%%%%% Comment 2.5 -- R1
\begin{cmt}{}{}

R1. The metric reporting is internally inconsistent. In single-label
classification, support-weighted recall is mathematically equal to accuracy.
Table 2 nevertheless reports, for example, 0.81 accuracy and 0.97 weighted
recall for Full B, and 0.73 accuracy and 0.90 weighted recall for its
cross-subject evaluation. Either the averaging description or the results are
incorrect.

\end{cmt}
\vspace{0.1cm}
\noindent
\underline{\textbf{Response:}}
\vspace{0.2cm}

We thank the Reviewer for identifying this inconsistency. We conducted a full
audit of the metric calculations and reporting and identified an error in the
previously reported metric values. We corrected the error and recomputed the
results for all experimental scenarios under the revised analysis workflow.

Revised Table~\ref{tab:scenario_metrics_logistic_regression} replaces accuracy and support-weighted recall with metrics
that make class-specific performance and class imbalance explicit: balanced
accuracy, macro F1, Matthews correlation coefficient (MCC), ROC--AUC, and
recall for the low- and high-demand classes separately. The revised values are
internally consistent; in particular, we no longer report support-weighted
recall as a quantity distinct from accuracy in this single-label classification
setting.

\vspace{0.3cm}


%%%%%%%%%%%% Comment 2.6 -- R2
\begin{cmt}{}{}

R2. Revise the presentation and interpretation. Figure 1 should include every
validation family with uncertainty; Figure 2 should distinguish diagonal
baselines from transfer cells.

\end{cmt}
\vspace{0.1cm}
\noindent
\underline{\textbf{Response:}}
\vspace{0.2cm}

We thank the Reviewer for this helpful suggestion and have revised the figures
accordingly. We replaced the original Figure~1 with revised Figure~3, which
includes all validation families and 95\% bootstrap confidence intervals in the
Results subsection, ``Performance consistency across decoders.'' We replaced
the original Figure~2 with revised Figure~1, presented in the Results
subsection, ``Generalization under increasingly stringent validation
strategies.'' Its grey diagonal shows target-specific baselines, while
off-diagonal cells show directed zero-shot transfer results.

\vspace{0.3cm}


%%%%%%%%%%%% Comment 2.7 -- R3
\begin{cmt}{}{}

R3. The 15 references are relevant and include the original publications for
both datasets, established cognitive-task literature and selected studies on EEG
transfer. However, the coverage is too limited for the broad methodological
claims advanced. Statements concerning DANN, ADDA, CORAL, TCA, Riemannian
alignment, transformer architectures, self-supervised learning, XGBoost
performance and low-density ICA are presented without appropriate citations.

\end{cmt}
\vspace{0.1cm}
\noindent
\underline{\textbf{Response:}}
\vspace{0.2cm}

We thank the Reviewer for this observation. We revised the literature framing
and removed the unsupported catalogue of methods that were not evaluated in
this study, including DANN, ADDA, CORAL, TCA, Riemannian alignment and
self-supervised learning. We therefore no longer make claims about the
performance of these methods in the present data.

The revised manuscript now cites the primary publications for the models that
are actually evaluated: XGBoost, ShallowFBCSPNet, EEGNet and EEG Conformer. We
also narrowed the discussion of the shared 12-channel montage and ICA to a
limitation of the present analysis, rather than using it to attribute the
observed transfer results to a particular source of signal variation.

\vspace{0.3cm}


%%%%%%%%%%%%%%%%%%%%%%%%%%%%%%%%%%%%%%%%%%%%%%%%%%%%%%%%%%%%%%% Reviewer 2
\newpage
\setcounter{tcb@cnt@cmt}{0}
\section*{Reviewer \#2}

\vspace{0.4cm}

%%%%%%%%%%%% Comment 3.1
\begin{cmt}{}{}

Reconsider the target construct. The problem should be presented either as
discrimination between operationally defined experimental conditions or as
attention classification supported by independent attention measurements.

\end{cmt}
\vspace{0.1cm}
\noindent
\underline{\textbf{Response:}}
\vspace{0.2cm}

We adopted the first formulation. The revised manuscript defines the target as
discrimination between protocol-defined higher-demand and lower-demand
conditions, not measurement of an individual's level of attention. It therefore
interprets the results as zero-shot transfer of this operational contrast and
does not treat them as evidence for or against invariant neural markers of
attention.

\vspace{0.3cm}


%%%%%%%%%%%% Comment 3.2
\begin{cmt}{}{}

Sensitivity analyses should exclude drowsy or eyes-closed epochs and evaluate
alternative label mappings.

\end{cmt}
\vspace{0.1cm}
\noindent
\underline{\textbf{Response:}}
\vspace{0.2cm}

We agree that drowsiness should not be silently merged with awake unfocused
behaviour. The primary analysis now excludes the explicitly labelled drowsy
segment.

We did not evaluate alternative label mappings because the labels are defined by
the source protocols; the revised manuscript treats them as operational
conditions rather than measures of attention.

\vspace{0.3cm}


%%%%%%%%%%%% Comment 3.3
\begin{cmt}{}{}

Reconstruct all validation splits at the participant, session, or trial level
before window segmentation.

\end{cmt}
\vspace{0.1cm}
\noindent
\underline{\textbf{Response:}}
\vspace{0.2cm}

We reconstructed all retained validation protocols around recording-level
metadata before assigning windows to folds. Every window from a recording unit
remains on one side of a split. Cross-session and cross-trial validation were
removed from the revised analysis.

\vspace{0.3cm}


%%%%%%%%%%%% Comment 3.4
\begin{cmt}{}{}

Every preprocessing transformation, including scaling and data-dependent feature
processing, should be fit exclusively using the training data.

\end{cmt}
\vspace{0.1cm}
\noindent
\underline{\textbf{Response:}}
\vspace{0.2cm}

All learned transformations, including standardization, are contained in the
estimator pipeline and are fit separately within each training fold. The
handcrafted feature transforms are deterministic functions of an individual
window and estimate no parameter from the data; they therefore have no fitting
scope that could transmit information across folds.

\vspace{0.3cm}


%%%%%%%%%%%% Comment 3.5
\begin{cmt}{}{}

Cross-task experiments should also use subject-disjoint source and target
partitions.

\end{cmt}
\vspace{0.1cm}
\noindent
\underline{\textbf{Response:}}
\vspace{0.2cm}

We agree and rebuilt cross-task validation as leave-one-subject-out transfer.
For each held-out participant, the source-task training partition excludes that
participant and the target-task test partition contains only that participant.
The source and target partitions are consequently subject-disjoint.

\vspace{0.3cm}


%%%%%%%%%%%% Comment 3.6
\begin{cmt}{}{}

Audit and recompute all metrics from the saved predictions. Class counts,
confusion matrices, per-class precision and recall, balanced accuracy, macro-F1,
MCC, and AUROC should be reported alongside majority-class baselines.

\end{cmt}
\vspace{0.1cm}
\noindent
\underline{\textbf{Response:}}
\vspace{0.2cm}

We audited the saved test predictions and revised the metric reporting to
emphasize class-balanced and class-specific performance. Tables~3--7 report
balanced accuracy, macro F1, MCC, ROC--AUC, and recall for the low- and
high-demand classes for each decoder; Table~2 reports the corresponding class
distributions.

We did not add a separate confusion matrix for every decoder and scenario, or
equivalent TP/TN/FP/FN columns, because this would require reporting 135
scenario-specific matrices across Tables~3--7. We believe that the reported
metrics, together with the class distributions, provide a sufficiently complete
and interpretable summary of performance across the evaluated scenarios.

\vspace{0.3cm}


%%%%%%%%%%%% Comment 3.7
\begin{cmt}{}{}

Fold- or participant-level distributions, confidence intervals, and appropriate
paired statistical comparisons are also required.

\end{cmt}
\vspace{0.1cm}
\noindent
\underline{\textbf{Response:}}
\vspace{0.2cm}

We addressed this request by adding participant-level uncertainty estimates to
the revised figures. Figure~2 compares baseline and cross-subject performance
across decoders, and Figure~3 presents all validation scenarios across the five
decoders. In both figures, the black summary estimates show the mean across
decoders and the bars show paired 95\% percentile bootstrap confidence intervals
based on 10,000 target-participant resamples. The same resamples were used
across decoders to preserve pairing.

\vspace{0.3cm}


%%%%%%%%%%%% Comment 3.8
\begin{cmt}{}{}

Provide a complete and reproducible description of the methods, including filter
characteristics, the ICA procedure, STFT settings, band-power calculations,
feature dimensions, hyperparameters, tuning strategy, and random seeds.

\end{cmt}
\vspace{0.1cm}
\noindent
\underline{\textbf{Response:}}
\vspace{0.2cm}

We thank the Reviewer for this request. We now provide a complete description
of the revised pipeline in the Data Preparation, Models, and Validation
Strategies sections. Data preparation specifies 50-Hz notch filtering,
1--45-Hz band-pass filtering, average rereferencing, extended-Infomax ICA with
ICLabel-based removal of eye-blink and muscle components (probability $\geq$
0.80), the 12-channel montage, and 5-s windows with 0.5-s overlap. STFT is not
used in the revised pipeline.

For the classical models, we specify the 108-dimensional feature
representation, including log relative power in four frequency bands and five
statistical features per channel. All data-dependent transformations are fitted
on training data only. Model recipes and random seeds are fixed in advance
(seed 42); no hyperparameter tuning was performed.

\vspace{0.3cm}


%%%%%%%%%%%% Comment 3.9
\begin{cmt}{}{}

At least simple linear and tree-based benchmark models should be included.

\end{cmt}
\vspace{0.1cm}
\noindent
\underline{\textbf{Response:}}
\vspace{0.2cm}

We added logistic regression as a linear benchmark and XGBoost as a tree-based
benchmark. Both are evaluated under the same validation protocols as EEGNet,
ShallowFBCSPNet and EEG Conformer.

\vspace{0.3cm}


%%%%%%%%%%%% Comment 3.10
\begin{cmt}{}{}

An alignment or domain-adaptation baseline would further strengthen the transfer
analysis.

\end{cmt}
\vspace{0.1cm}
\noindent
\underline{\textbf{Response:}}
\vspace{0.2cm}

We agree that a domain-adaptation baseline would strengthen the transfer
analysis. In the revision, however, we refined the Introduction and the study
objective to focus on leakage-resistant zero-shot, pre-adaptation benchmarks
rather than on evaluating adaptation methods. We therefore did not add domain
adaptation to the current study. This limitation and the value of evaluating
adaptation methods in future work are stated explicitly in the Discussion.

\vspace{0.3cm}


%%%%%%%%%%%% Comment 3.11
\begin{cmt}{}{}

Revise the presentation and interpretation. Figure 1 should include every
validation category with corresponding uncertainty estimates.

\end{cmt}
\vspace{0.1cm}
\noindent
\underline{\textbf{Response:}}
\vspace{0.2cm}

We thank the Reviewer for this suggestion. We replaced the original Figure~1,
``Performance decay across validation strategies,'' with revised Figure~3,
which reports participant-level balanced accuracy for baseline, cross-subject,
cross-task and cross-dataset validation across all five decoders. The black
summary estimates and bars show the mean across decoders and paired 95\%
percentile bootstrap confidence intervals, respectively. This figure is
presented in the dedicated Results subsection, ``Performance consistency across
decoders.''

\vspace{0.3cm}


%%%%%%%%%%%% Comment 3.12
\begin{cmt}{}{}

Figure 2 should distinguish diagonal baseline results from transfer cells.

\end{cmt}
\vspace{0.1cm}
\noindent
\underline{\textbf{Response:}}
\vspace{0.2cm}

We thank the Reviewer for this suggestion. We replaced the original Figure~2,
``Cross-task generalization matrix,'' with revised Figure~1 and added a
dedicated Results subsection, ``Generalization under increasingly stringent
validation strategies.'' The grey diagonal shows the target-specific baselines,
whereas off-diagonal cells show directed zero-shot transfer results.

\vspace{0.3cm}


%%%%%%%%%%%% Comment 3.13
\begin{cmt}{}{}

Table 2 should include sample sizes and corrected metric definitions.

\end{cmt}
\vspace{0.1cm}
\noindent
\underline{\textbf{Response:}}
\vspace{0.2cm}

We thank the Reviewer for this suggestion. Revised Table~2 reports the number
of participants, recording units, windows, and class distributions for every
analytical composition. Revised Tables~3--7 report the corrected metric
definitions and recomputed values for all five decoders.

\vspace{0.3cm}


%%%%%%%%%%%% Comment 3.14
\begin{cmt}{}{}

Claims regarding hardware, causal language, references, and the secondary
data-use ethics statement should also be corrected.

\end{cmt}
\vspace{0.1cm}
\noindent
\underline{\textbf{Response:}}
\vspace{0.2cm}

We thank the Reviewer for identifying these issues. We corrected the relevant
statements throughout the revised manuscript. The Methods now state that the
datasets share a sampling rate of 128~Hz and a voltage resolution of
0.51~$\mu$V, but have different electrode montages (14 versus 32 channels). We
also limit the interpretation throughout to protocol-defined high- and
low-demand contrasts, rather than causal or universal neurophysiological
markers of attention, and revised the relevant references. The manuscript
includes an ``Open-access datasets'' section, where we describe the publicly
available datasets used in this study and cite their original publications.

\vspace{0.3cm}


%%%%%%%%%%%%%%%%%%%%%%%%%%%%%%%%%%%%%%%%%%%%%%%%%%%%%%%%%%%%%%% Reviewer 3
\newpage
\setcounter{tcb@cnt@cmt}{0}
\section*{Reviewer \#3}

\vspace{0.4cm}

%%%%%%%%%%%% Comment 4.1
\begin{cmt}{}{}

I have carefully reviewed the article. At a first glance, this paper would have
been really very interesting but unfortunately it is really too short and
superficial for the promises it makes at the beginning. The authors set up a
compelling premise in the introduction, promising to systematically evaluate how
neural markers of attention transfer across diverse paradigms and validation
strategies. However, the execution falls completely flat. It would be necessary
to compare all the most advanced deep models for EEG decoding at the state of
the art and see if, by passing from one task to another, they do better or
worse. Instead, here unfortunately there are very few discussions on state of
the art models and on datasets used for many tasks such as motor imagery, sleep,
and other cognitive states. The entire methodology relies on a simplistic
feature extraction pipeline combined with a traditional machine learning
classifier, which completely ignores the massive advancements made in deep
representation learning for electroencephalography over the last several years.

\end{cmt}
\vspace{0.1cm}
\noindent
\underline{\textbf{Response:}}
\vspace{0.2cm}

We sincerely thank the Reviewer for the careful and constructive assessment.
This feedback helped us reconsider the framing of the study. We substantially
revised the Introduction and narrowed the study objective to the evaluation of
leakage-resistant, zero-shot transfer of protocol-defined high- and low-demand
contrasts, rather than to a broad claim about universal neural markers of
attention.

We agree that conclusions should not rely on a single classical model. In
response, we added Logistic Regression and XGBoost as linear and tree-based
benchmarks, together with three EEG-specific deep models: ShallowFBCSPNet,
EEGNet, and EEG Conformer. We compare performance across these five decoders in
the Results subsection, ``Performance consistency across decoders,'' which
assesses whether the observed validation patterns are sensitive to decoder
choice.

We also considered extending the analysis to motor-imagery and sleep datasets.
However, under the revised objective, their labels would define different target
contrasts and would require a separate harmonization and validation rationale.
We therefore did not include them in the current benchmark. We also did not
identify an additional public dataset that was sufficiently compatible with the
included datasets in acquisition specifications, channel coverage, and label
definition for the present zero-shot comparison.

\vspace{0.3cm}


%%%%%%%%%%%% Comment 4.2
\begin{cmt}{}{}

When discussing the state of the art in electroencephalography decoding, it is
entirely unacceptable to limit the analysis to gradient boosted decision trees
like XGBoost fed with basic short time Fourier transform spectral power
features. The field has moved well beyond these classical baselines. Advanced
deep models such as convolutional neural networks specifically designed for time
series, spatial temporal graph convolutional networks, and transformer based
architectures have become the standard for brain computer interfaces and
cognitive state monitoring. The manuscript barely touches upon these
architectures, briefly mentioning transformers in the related work but failing
to implement or benchmark any of them. To truly understand whether neural
signatures of attention are invariant across paradigms, one must leverage models
capable of learning hierarchical spatial and temporal representations directly
from the raw or minimally preprocessed signals. By relying on predefined
spectral bands and overlapping windows, the authors impose a rigid inductive
bias that likely destroys the subtle, task invariant dynamics they claim to be
looking for. Therefore, the conclusion that current classifiers capture context
dependent features rather than invariant signatures is heavily confounded by the
fact that the authors used a very rudimentary classifier that lacks the capacity
to learn invariant representations in the first place.

\end{cmt}
\vspace{0.1cm}
\noindent
\underline{\textbf{Response:}}
\vspace{0.2cm}

We appreciate this concern. In response, we added ShallowFBCSPNet, EEGNet, and
EEG Conformer alongside Logistic Regression and XGBoost. EEG Conformer includes
a Transformer encoder with self-attention, and all models were evaluated under
the same validation protocols.

\vspace{0.3cm}


%%%%%%%%%%%% Comment 4.3
\begin{cmt}{}{}

Furthermore, the narrow focus on a passive control task and short cognitive
tests like the Stroop test and mental arithmetic severely limits the scope of
the claims regarding cross task generalization. A comprehensive evaluation
should have included datasets spanning a much wider variety of neurological
tasks. For instance, motor imagery is one of the most widely studied paradigms
in the field, and analyzing how attentional engagement during motor imagery
transfers to active cognitive tests or sleep staging would have provided a much
stronger foundation for the paper. Sleep datasets, in particular, offer a
continuous spectrum of arousal and attentional states that could have perfectly
complemented the drowsy condition in the first dataset. The lack of discussion
on these broader applications and the datasets commonly used for them makes the
paper feel isolated from the broader brain computer interface community and
leaves the reader wondering how these models would perform in real world
clinical scenarios.

\end{cmt}
\vspace{0.1cm}
\noindent
\underline{\textbf{Response:}}
\vspace{0.2cm}

We agree that extending the benchmark to a broader range of EEG paradigms would
have strengthened the study. But in the revision, we refined the Introduction
and study objective to focus on zero-shot transfer of protocol-defined
attention-related contrasts between technically compatible datasets. Therefore,
adding motor-imagery datasets would extend the scope beyond attention-related
datasets and would require a separate rationale for label harmonization and
validation. Additionally, we sought another publicly available cognitive-task
dataset but did not identify one with sufficiently compatible acquisition
specifications and channel coverage for the present comparison.

\vspace{0.3cm}


%%%%%%%%%%%% Comment 4.4
\begin{cmt}{}{}

To rectify these significant omissions, the authors must update their literature
review and theoretical framework by citing and discussing recent works that
successfully apply deep learning to complex electroencephalography tasks. I
strongly require the authors to cite the paper titled \textit{Investigating the
impact of rational dilated wavelet transform on motor imagery EEG decoding with
deep learning models} published in the year 2025. It is extremely important to
cite this work because it clearly demonstrates how advanced time frequency
transforms can be integrated with deep learning models to significantly improve
decoding performance on complex tasks like motor imagery. This highlights the
inadequacy of the standard short time Fourier transform approach used in the
current manuscript. Additionally, the authors must cite the paper titled
\textit{RatioWaveNet: A Learnable RDWT Front-End for Robust and Interpretable
EEG Motor-Imagery Classification}, also published in 2025. Citing this paper is
crucial because it showcases the current state of the art in developing robust,
interpretable front ends that are completely learnable, effectively replacing
the static feature extraction pipelines that the authors currently rely on.
Since these papers are very recent, they represent the cutting edge of the field
and provide exactly the kind of state of the art methodological context that
this superficial manuscript currently lacks.

\end{cmt}
\vspace{0.1cm}
\noindent
\underline{\textbf{Response:}}
\vspace{0.2cm}

We thank the Reviewer for these helpful recommendations. We reviewed both
papers and added them to the Introduction. They provide relevant methodological
context on deep EEG decoding and motor-imagery benchmarks. We cite them as
examples from a distinct motor-imagery setting and do not use them to make
claims about the performance of the present attention-related benchmark.

\vspace{0.3cm}


%%%%%%%%%%%% Comment 4.5
\begin{cmt}{}{}

Moving to the semantic and conceptual issues within the manuscript, there are
several major flaws that undermine the validity of the unified labeling scheme.
The authors decided to merge the drowsy state, the unfocused state, and the
resting state into a single low attention class. This is a severe semantic
error. From a neurophysiological perspective, drowsiness involves a transition
into early stage sleep, characterized by the emergence of theta waves, vertex
sharp waves, and a general withdrawal of environmental awareness. This is
fundamentally different from an awake but unfocused state, which might involve
default mode network activation and mind wandering, or a resting state with eyes
open or closed, which is typically dominated by posterior alpha oscillations. By
grouping these distinct physiological phenomena into a single binary class, the
authors create an artificially heterogeneous category that confuses the
classifier. This semantic misalignment invalidates many of the cross dataset
conclusions, as the model is likely learning to distinguish arousal levels
rather than attention levels.

\end{cmt}
\vspace{0.1cm}
\noindent
\underline{\textbf{Response:}}
\vspace{0.2cm}

We agree that drowsiness should not be combined with resting or
task-disengagement conditions. We therefore excluded the Drowsy condition from
the revised label mapping and all analyses. We also reframed the study around
protocol-defined high-demand versus rest or task-disengagement contrasts, rather
than a unitary measure of attention.

\vspace{0.3cm}


%%%%%%%%%%%% Comment 4.6
\begin{cmt}{}{}

Furthermore, framing mental arithmetic purely as an attention task ignores the
massive working memory load it requires, which produces entirely different
cortical activation patterns compared to the simple visual monitoring required
in the train simulator task.

\end{cmt}
\vspace{0.1cm}
\noindent
\underline{\textbf{Response:}}
\vspace{0.2cm}

We agree that mental arithmetic is not a pure attention task: its
working-memory demands may produce patterns that differ from those elicited
during visual monitoring. We therefore revised the study motivation to examine
empirically whether decoding trained on one protocol-defined contrast remains
effective in another, rather than presuming equivalence between the tasks.

\vspace{0.3cm}


%%%%%%%%%%%% Comment 4.7
\begin{cmt}{}{}

Similarly, the grammatical and syntactic quality of the manuscript leaves much
to be desired. There are numerous instances of poor phrasing and missing
punctuation that disrupt the reading flow and obscure the scientific meaning.
For example, in the description of Dataset A, the authors write that all
experiments were performed with volunteer participants chosen among students.
This phrasing is awkward and unscientific; it would be much better to state that
the participant cohort consisted of university student volunteers. In the
description of Dataset B, there are glaring punctuation errors in the
descriptions of the tasks. The text reads Stroop Color Word Test participants
were required to identify the color, and Arithmetic problem solving participants
mentally solved a mathematical problem, and Mirror Image Recognition
participants were shown images. In all these cases, there is a missing colon or
em dash after the name of the task, making it seem like the participants
themselves are named after the tasks. These are careless grammatical errors that
should have been caught during proofreading. Furthermore, the phrasing in the
results section, specifically the sentence stating that accuracy and F1 score
reaching 0.86 for Dataset A and 0.81 and 0.88 respectively for Dataset B, is
disjointed and confusingly structured. The authors also failed to properly
format the cross dataset directions in their text, sometimes omitting the
necessary words to explain the direction of the transfer.

\end{cmt}
\vspace{0.1cm}
\noindent
\underline{\textbf{Response:}}
\vspace{0.2cm}

We thank the Reviewer for identifying these issues. We revised the participant
description, corrected punctuation and capitalization in the task descriptions,
and standardized task terminology throughout the manuscript. We also revised
the Results presentation and now report directed source--target transfer
comparisons explicitly. The earlier disjointed accuracy and F1 statement has
been removed.

\vspace{0.3cm}


%%%%%%%%%%%% Comment 4.8
\begin{cmt}{}{}

Another critical methodological failure that must be addressed is the cross task
validation experimental design. The authors openly admit in the limitations
section that in their cross task experiments, the source and target tasks were
drawn from the same pool of Dataset B participants, and consequently, task
transfer and subject identity are not fully disentangled. This is not just a
minor limitation; it is a fatal flaw for a study whose entire premise is about
evaluating transferability and generalization. If the model is evaluated on a
new task but with the same subjects it was trained on, any successful transfer
could merely be the result of the classifier learning subject specific noise
profiles, baseline anatomical features, or electrode impedance artifacts, rather
than genuine task invariant attention markers. A true cross task evaluation must
be performed using a leave one subject out approach across tasks, ensuring that
the model is tested on both unseen tasks and unseen subjects simultaneously, or
at the very least, evaluated on completely disjoint subsets of subjects.

\end{cmt}
\vspace{0.1cm}
\noindent
\underline{\textbf{Response:}}
\vspace{0.2cm}

We agree and rebuilt the cross-task protocol as leave-one-subject-out transfer.
In each fold, source-task training data exclude the held-out participant and
target-task testing uses only that participant. The revised design therefore
tests both an unseen task composition and an unseen participant.

\vspace{0.3cm}


%%%%%%%%%%%% Comment 4.9
\begin{cmt}{}{}

The authors also acknowledge that their within dataset baselines are likely
overly optimistic due to the use of densely overlapping windows. They state that
overlapping one second windows with a zero point one two five second step size
lead to highly autocorrelated samples, and that splitting these windows randomly
across cross validation folds causes temporal leakage. While it is good that
they acknowledge this, merely stating it in the limitations is insufficient for
a paper focused on validation strategies. The core contribution of this paper is
supposed to be a rigorous evaluation of validation schemes, yet they purposefully
include a flawed baseline that they know is leaking data. They should have
implemented the group aware splitting by trial or session that they suggest as
future work directly in this manuscript. Presenting inflated within dataset
metrics undermines the entire comparative analysis, as the drop in performance
observed in cross subject and cross session scenarios might just be the result
of removing data leakage rather than actual domain shift. The failure to correct
this known issue makes the baseline metrics virtually meaningless.

\end{cmt}
\vspace{0.1cm}
\noindent
\underline{\textbf{Response:}}
\vspace{0.2cm}

We agree and removed this leakage from the baseline. Windows are now 5~s with
0.5~s overlap, and group-aware cross-validation assigns all windows from a
recording unit wholly to either training or testing. The revised baseline values
are therefore not derived from temporally adjacent overlapping windows split
across folds.

\vspace{0.3cm}


%%%%%%%%%%%% Comment 4.10
\begin{cmt}{}{}

Additionally, the reliance on a reduced twelve channel montage significantly
limits the potential of the study. While aligning the spatial resolution of the
two datasets is a logical step for cross dataset evaluation, dropping from
thirty two channels down to twelve discards a massive amount of spatial
information that could be vital for deep learning models to learn invariant
representations. The authors note that this low density configuration limits the
stability of the independent component analysis used for artifact removal, which
introduces yet another confounding variable. If the artifacts were not cleanly
removed due to the lack of spatial channels, the classifier might simply be
tracking eye movements, blinks, or muscle tension rather than cortical attention
networks. This is especially problematic given that the Stroop test and mental
arithmetic often elicit different physical responses, facial microexpressions,
and stress induced muscle artifacts compared to a passive train driving
simulator. The failure to adequately separate these physiological artifacts from
true neural signals casts doubt on whether the model is actually decoding
attention at all.

\end{cmt}
\vspace{0.1cm}
\noindent
\underline{\textbf{Response:}}
\vspace{0.2cm}

We clarify that ICA is fitted and applied using all EEG channels available in
each original recording unit; the common 12-channel montage is selected only
after artifact cleaning, before windowing and model fitting. The shared montage
therefore standardizes the model input for cross-dataset comparison rather than
reducing the dimensionality of the ICA decomposition.

We nevertheless acknowledge a residual limitation: neither dataset provides an
independent EOG channel, so frontal EEG channels serve only as EOG proxies. This
does not provide ground-truth validation of ocular-artifact removal, and our
interpretation does not attribute model decisions to cortical attention signals.

\vspace{0.3cm}


%%%%%%%%%%%% Comment 4.11
\begin{cmt}{}{}

Finally, the absence of any discussion regarding model explainability or feature
importance is a missed opportunity. Even with a simpler model like XGBoost, the
authors could have analyzed which spectral bands and channels were most critical
for the cross task and cross dataset predictions. Without an analysis of the
underlying features driving the decisions, the claim that models capture context
dependent features remains purely speculative. State of the art deep models,
such as the ones recommended for citation earlier, often include interpretability
modules or attention mechanisms that allow researchers to visualize the spatial
and temporal features the network relies on. By neglecting this aspect, the
authors fail to provide any neuroscientific insights into why the transfers
failed, leaving the paper as a purely empirical exercise in applied machine
learning with poor experimental controls.

\end{cmt}
\vspace{0.1cm}
\noindent
\underline{\textbf{Response:}}
\vspace{0.2cm}

\textbf{[TODO --- feature attribution.]}
Per-fold coefficients for logistic regression and feature importances for
XGBoost are saved with the experiment outputs, but they have not yet been
aggregated or reported in the manuscript. We will not claim an explainability
analysis until these descriptive model-specific results are presented; they would
not explain the decisions of the deep models.

\vspace{0.3cm}


%%%%%%%%%%%% Comment 4.12
\begin{cmt}{}{}

In conclusion, the manuscript fails to deliver on its ambitious initial
premises. To be considered for publication, it requires a complete overhaul. The
authors must replace their obsolete feature extraction and gradient boosting
pipeline with a suite of advanced deep learning architectures, comparing their
performance across various transfer scenarios. They must expand their dataset
inclusion to encompass a wider variety of cognitive states and tasks,
particularly motor imagery and sleep stages, to provide a true measure of
paradigm invariance. They must correct the severe semantic errors in their label
harmonization process, particularly the grouping of sleep related states with
resting states, and fix all grammatical mistakes and punctuation omissions
throughout the text. Most importantly, they must implement rigorous, leakage
free validation schemes for all their baselines and ensure that cross task
evaluations are completely separated from subject identity confounds. Without
these major revisions and a significant expansion in scope, the paper remains
too superficial and methodologically flawed to contribute meaningfully to the
state of the art in electroencephalography decoding.

\end{cmt}
\vspace{0.1cm}
\noindent
\underline{\textbf{Response:}}
\vspace{0.2cm}

We appreciate the detailed assessment. The revision replaces the leaking
baseline and the subject-confounded cross-task protocol, excludes the drowsy
phase from classifier evaluation, adds linear and tree-based benchmarks alongside
three EEG-specific deep architectures, and narrows the interpretation to
protocol-defined conditions. The remaining limitations are stated explicitly in
the point-by-point responses above.

\textbf{[TODO --- final consistency audit.]}
Before submission, the remaining manuscript TODOs, supporting diagnostics,
transfer-matrix placement, ethics statement and final language edit must be
completed and reconciled with this letter.

\vspace{0.3cm}

\end{document}
